\documentclass[runningheads]{llncs}

\usepackage[T1]{fontenc}
\usepackage{graphicx}
\usepackage{booktabs}
\usepackage{hyperref}
\usepackage{amsmath}
\usepackage{array}
\usepackage{makecell}
\usepackage{multirow}
\usepackage[utf8]{inputenc}
\usepackage[final]{microtype}
\usepackage{float}
\setlength{\emergencystretch}{3em}
\setlength{\floatsep}{8pt plus 2pt minus 2pt}
\setlength{\textfloatsep}{10pt plus 2pt minus 4pt}
\setlength{\intextsep}{8pt plus 2pt minus 2pt}

\hypersetup{colorlinks=true,linkcolor=blue,urlcolor=blue,citecolor=blue}

\begin{document}

\title{In-Context Learning: Chain-of-Thought vs.\ Tree-of-Thoughts
Prompting on Multi-Step Reasoning}

\titlerunning{ICL: CoT vs.\ ToT on Multi-Step Reasoning}

\author{Sercan Utku Temelci \and
O\u{g}uzhan \c{C}akal}

\authorrunning{S.~Utku Temelci and O\u{g}uzhan \c{C}akal}

\institute{Department of Computer Engineering,
\'{I}zmir Institute of Technology (\'{I}YTE), \'{I}zmir, Turkey\\
\email{\{sercanutkutemelci, oguzhancakal\}@std.iyte.edu.tr}}

\maketitle

\begin{abstract}
We present an empirical comparison of four In-Context Learning (ICL)
prompting strategies on the GSM8K grade-school mathematics benchmark:
Zero-Shot Chain-of-Thought (CoT), Few-Shot CoT (8-shot),
Tree-of-Thoughts (ToT) with BFS exploration at two search depths,
and Self-Consistency voting.
All experiments use Llama~3.1~8B via the Groq API without fine-tuning.
Zero-Shot CoT achieves \textbf{95.0\%} accuracy, Self-Consistency
\textbf{92.0\%}, and ToT (BFS-2, depth=6) together with Few-Shot CoT
each reach \textbf{85.0\%}---all on the same 100-problem seed-42 sample.
After an answer-extraction format fix, Zero-Shot CoT alone sits on the
Pareto frontier of accuracy and inference cost at the 8B scale.
A depth ablation (depth=3 vs.\ depth=6) yields only $+4.0$~pp at
$2.1\times$ inference cost, indicating diminishing returns when the
evaluator is the same 8B backbone.
A cross-model replication on Claude Haiku~4.5 with full evaluation
identity preserved collapses the inter-strategy gap from $14$~pp to
$2$~pp, revealing a \emph{ceiling effect}: deliberation- and
ensemble-style strategies pay off only when the base model has
headroom to lose and identifiable wrong branches to prune.
A pairwise Cohen's $\kappa$ analysis confirms this picture
independently---ToT depths on Llama reach near-zero agreement
($\kappa{=}0.011$) versus perfect agreement on Haiku
($\kappa{=}1.000$).
\end{abstract}

\keywords{In-Context Learning \and Chain-of-Thought \and
Tree-of-Thoughts \and GSM8K \and Large Language Models}

%------------------------------------------------------------------------
\section{Introduction}
\label{sec:intro}

Large Language Models (LLMs) have demonstrated remarkable capabilities
in natural language tasks through In-Context Learning (ICL), where the
model adapts to new tasks from demonstrations provided in the prompt
without updating model weights~\cite{brown2020gpt3}.
This ability to generalise from a handful of examples has made ICL
the dominant paradigm for practical LLM deployment, motivating active
research into how the structure and content of prompts influence
multi-step reasoning quality.
Among ICL techniques, Chain-of-Thought (CoT) prompting has emerged as
a particularly effective approach, guiding models to produce
intermediate reasoning steps before arriving at a final
answer~\cite{wei2022cot}.
The zero-shot variant~\cite{kojima2022zeroshot} further simplified
this by showing that a single trigger phrase suffices for
instruction-tuned models to produce step-by-step reasoning, suggesting
that structured reasoning has been internalised during pre-training.

Tree-of-Thoughts (ToT) extends this paradigm by exploring multiple
reasoning branches simultaneously and pruning low-quality paths through
a separate evaluation step~\cite{yao2023tot}.
While ToT has shown impressive gains over CoT on GPT-4, its
effectiveness on smaller, open-weight models remains unexplored.
A central open question is whether the multi-branch search ToT
provides translates to accuracy gains when the evaluator model is the
same small backbone---or whether evaluator weakness undermines the
pruning mechanism entirely.

This paper presents a controlled empirical comparison of four
prompting strategies (Zero-Shot CoT, Few-Shot CoT (8-shot), ToT
with BFS, and Self-Consistency) on the GSM8K
benchmark~\cite{cobbe2021gsm8k}, using Llama~3.1~8B as a common
backbone.
Zero-Shot CoT reaches 95.0\% on 100 samples (after an
answer-extraction format fix, Section~\ref{sec:setup}), our ToT
implementation (depth=6, BFS with branching factor 2) reaches
85.0\%, and Self-Consistency voting reaches 92.0\%.
A depth ablation (\texttt{depth}=3 vs.\ \texttt{depth}=6) further shows
that doubling the search depth yields only $+4.0$\,pp accuracy at
$2.1\times$ the inference cost, suggesting that the bottleneck for
ToT at the 8B scale is evaluator discrimination rather than search
depth.
To test whether these findings are model-specific, we replicate all
five configurations on Claude Haiku~4.5, a frontier-tier model, with
exact sample, prompt, and evaluation identity preserved---a controlled
comparison that reveals a scale-dependent ceiling effect.

%------------------------------------------------------------------------
\section{Related Work}
\label{sec:related}

The foundational prompting strategy for multi-step reasoning is
Chain-of-Thought (CoT), introduced by Wei et al.~\cite{wei2022cot},
who demonstrated that providing few-shot examples with step-by-step
solutions substantially improves LLM performance on arithmetic and
symbolic reasoning.
Kojima et al.~\cite{kojima2022zeroshot} subsequently showed that
simply appending \textit{``Let's think step by step''} to a query
elicits CoT reasoning in a zero-shot setting, establishing a
lightweight baseline that avoids hand-crafted demonstrations and
anchors our Zero-Shot CoT configuration.

Building directly on CoT, Yao et al.~\cite{yao2023tot} proposed
Tree-of-Thoughts (ToT), which frames reasoning as a search problem
over a tree of intermediate ``thoughts''.
At each node, the model generates multiple candidate continuations,
evaluates them with a separate scoring prompt (assigning
\textit{sure/maybe/impossible}), and retains only the most promising
branches using BFS or DFS.
Their experiments with GPT-4 showed substantial gains on tasks
requiring deliberate planning; however, all results used models with
100B+ parameters, leaving the effectiveness of the same mechanism on
smaller open-weight evaluators as an open question.

A complementary ensemble approach was proposed by Wang et
al.~\cite{wang2022selfconsistency}, who showed that sampling diverse
CoT paths at temperature $>0$ and selecting the final answer by
majority vote improves CoT accuracy by 10--20~pp without changing the
prompt.
Self-Consistency thus occupies a middle ground between the
minimal-cost single-chain CoT and the expensive multi-branch ToT.

All three strategies are instances of the In-Context Learning paradigm
established by Brown et al.~\cite{brown2020gpt3}: LLMs can generalise
to unseen tasks from a small number of in-context demonstrations
without gradient updates.
This shared grounding makes the four strategies directly comparable
under a fixed model and prompt template, as in our evaluation.

%------------------------------------------------------------------------
\section{Methodology}
\label{sec:method}

\subsection{Zero-Shot Chain-of-Thought}
The problem is presented with the trigger phrase
\textit{``Let's think step by step.''}~\cite{kojima2022zeroshot}
plus the output-format instruction
\textit{``End your response with: The answer is X.''}
to prevent trailing-modifier parse errors
(see Section~\ref{sec:setup}, \emph{Answer-extraction format fix}).

\subsection{Few-Shot Chain-of-Thought (8-shot)}
Eight worked examples from Wei et al.~\cite{wei2022cot} are prepended
to the query, each containing a full reasoning chain ending with
\textit{``The answer is X.''}\ followed by the test problem.

\subsection{Tree-of-Thoughts (BFS, branching=2, depth=6)}
\label{sec:method:totfixed}
ToT (\texttt{src/run\_tot\_fixed.py}) explores reasoning as a search over
a tree of intermediate ``thoughts''.
At each depth level, two candidate ``next-step thoughts'' are generated
with temperature~$= 0.7$.
A separate evaluator prompt scores each thought as
\textit{sure}, \textit{maybe}, or \textit{impossible}; only \textit{sure}
and \textit{maybe} thoughts are expanded in the next depth level.
After BFS completion, a dedicated extraction step queries the best leaf
node for the final numeric answer.
The configuration (branching factor 2, depth 6) requires approximately
45 API calls per problem.

\subsection{Self-Consistency (5 paths)}
Few-Shot CoT is sampled 5 times per problem with temperature~$= 0.7$.
The final answer is selected by majority vote over the 5 sampled
numeric answers~\cite{wang2022selfconsistency}.

\subsection{Cross-Model Evaluation Protocol}
\label{sec:method:crossmodel}
To test whether strategy orderings generalise beyond the 8B scale, we
re-run all five configurations on Claude Haiku~4.5 using the same
100-problem sample, prompts, parser, and gold-label correction.
Sample, prompt, and parser identity are preserved so that any accuracy
difference is attributable to the backbone alone.
The three single-call strategies are batched through the Anthropic
\texttt{messages.batches} API (50\% discount); ToT runs synchronously
because its BFS state cannot be safely flattened into batch items.
Realized cost and token counts are in Section~\ref{sec:results:crossmodel}.

%------------------------------------------------------------------------
\section{Experimental Setup}
\label{sec:setup}

\paragraph{Dataset.}
We use GSM8K~\cite{cobbe2021gsm8k}, a benchmark of 8,792 grade-school
math word problems requiring 2--8 chained reasoning steps.
All five configurations (Zero-Shot CoT, Few-Shot CoT, ToT depth=3,
ToT depth=6, and Self-Consistency) use the same 100-problem random
sample of GSM8K's 1,319-item test split (seed~$= 42$).
One item in our sample contains a gold-label annotation error
(see the \emph{GSM8K annotation error} paragraph below); we manually
verified the correct value and use the corrected label in all
evaluations, preserving the original sample sizes.

\paragraph{Model.}
All experiments use Llama~3.1~8B (\texttt{llama-3.1-8b-instant}) via
the Groq inference API.
No fine-tuning is applied; the model is used in its base instruction
configuration.

\paragraph{Evaluation.}
The primary metric is \textit{exact-match accuracy}: the predicted
final numeric answer is compared against the gold answer after stripping
commas and converting to float.
Answers are extracted with a regular expression that captures the
gold-format marker
\texttt{\#\#\#\#}; if no such marker is found, the last number in
the response is used.
We additionally report the \textit{parse failure rate}: the fraction
of responses from which no numeric prediction could be extracted.
Across all experiments, the \texttt{\#\#\#\#} gold-format marker was
never spontaneously produced by Llama~3.1~8B; every response fell back
to the last-number heuristic, which succeeded in all 500 evaluated
responses (parse failure rate~$= 0\%$).

\paragraph{Answer-extraction format fix.}
An initial Zero-Shot CoT run sometimes ended with a trailing-modifier
number rather than the substantive answer (e.g.\ ``\textit{36 hours
in 4 weeks}'' → parser captures \texttt{4} instead of \texttt{36}).
Appending ``\textit{End your response with: The answer is X.}'' to the
prompt raised Zero-Shot CoT from \textbf{88.0\%} to \textbf{95.0\%}
($+7$\,pp), concentrated in the hard bucket.
All Zero-Shot CoT numbers use this format-aligned prompt; the other
strategies already produce unambiguous final answers via their
demonstrations or dedicated extraction prompts.

\paragraph{GSM8K annotation error (gold-label correction).}
Manual inspection identified one annotation error in our sample (the
\textit{carnival} item). GSM8K's own calculator-annotated solution
string reveals the mechanism: \texttt{750+430+400+700=2280}, i.e.\ the
final sum substitutes an intermediate value from the previous step
(Maryam's \$400 excess over Sarah) for Sarah's own \$300 contribution,
which never appears as a term; the correct total is
$430+750+300+700=\$2{,}180$. We adopt the corrected label for all
reported accuracies. On Llama, Zero-Shot CoT and ToT depth=6 (both
produced 2{,}180) shift from incorrect to correct; on Haiku, all five
strategies produced 2{,}180 and each gains a uniform $+1$\,pp. In both
cases the correction leaves strategy ranking unchanged, consistent with
the $\sim$1--2\% annotation noise floor of GSM8K~\cite{cobbe2021gsm8k}.

\paragraph{Infrastructure.}
Rate-limit constraints on the Groq free tier required an 8-second
inter-call pause (CoT) and a 30-second inter-problem pause (ToT).
All random seeds are fixed at 42 for reproducibility.
All code and result files are publicly available.\footnote{\url{https://github.com/OguzhanCKL/CENG467_In-Context_Learning}}

\paragraph{Secondary model and cross-model setup.}
\label{sec:setup:crossmodel}
The secondary backbone is Claude Haiku~4.5
(\texttt{claude-haiku-4-5-20251001}) via the Anthropic Python SDK.
ZS-CoT, FS-CoT, and SC requests (700 total) are submitted via the
\texttt{messages.batches} API at a 50\% discount; ToT runs
synchronously (BFS state cannot be batched).
All Haiku runs use the same seed-42 sample, prompts, parser, and
carnival gold-label correction as the Llama runs.
Realized total cost: \$3.26 across all five strategies
(\$0.52 batch + \$0.72 ToT-d3 + \$2.02 ToT-d6).

%------------------------------------------------------------------------
\section{Results}
\label{sec:results}

\subsection{Main Results}

Table~\ref{tab:main} summarises the accuracy of the four strategies.

\begin{table}[H]
\centering
\caption{Main results on GSM8K. Llama~3.1~8B, Groq API.
Accuracies use the GSM8K gold labels with one item corrected (see
Section~\ref{sec:setup}, \emph{GSM8K annotation error}).
$\bar{T}$ denotes the average response length in tokens (one token
$\approx$ four characters), aggregated across all API calls per problem;
TPC = tokens per correct answer (lower is better).}
\label{tab:main}
\begin{tabular}{@{}lcccccl@{}}
\toprule
\textbf{Strategy} & \textbf{Accuracy} & \textbf{N} &
\textbf{API Calls} & \boldmath$\bar{T}$ & \textbf{TPC} & \textbf{Status} \\
\midrule
Zero-Shot CoT                     & \textbf{95.0\%} & 100 & $\sim$1  & 189  & 199  & Complete \\
Few-Shot CoT (8-shot)             & \textbf{85.0\%} & 100 & $\sim$1  & 145  & 171  & Complete \\
Tree-of-Thoughts (BFS-2, depth=3) & \textbf{81.0\%} & 100 & $\sim$21 & 315  & 389  & Complete (ablation) \\
Tree-of-Thoughts (BFS-2, depth=6) & \textbf{85.0\%} & 100 & $\sim$45 & 561  & 660  & Complete \\
Self-Consistency (5 paths)        & \textbf{92.0\%} & 100 & 5        & 738  & 802  & Complete \\
\bottomrule
\end{tabular}
\end{table}

Beyond accuracy, we report two efficiency metrics that have become
standard in recent reasoning benchmarks~\cite{ockbench2025,tale2024}:
the average per-problem token budget $\bar{T}$ (a model-agnostic proxy
for compute cost) and \textit{tokens per correct answer}
($\text{TPC} = \bar{T} / \text{accuracy}$), which captures the
practical cost of \emph{useful} reasoning---a method achieving high
accuracy but spending disproportionately many tokens may be
impractical~\cite{yao2023tot,ockbench2025}.
Few-Shot CoT is the most token-efficient strategy (171 TPC), followed
closely by Zero-Shot CoT (199); both ToT configurations spend
2--4$\times$ more tokens per correct answer (389 and 660), and
Self-Consistency---despite the highest accuracy---is the most expensive
at 802~TPC, consistent with its 5$\times$ sample budget.

Zero-Shot CoT achieves 95.0\% (after the answer-extraction
format fix, Section~\ref{sec:setup}) and substantially
outperforms Few-Shot CoT (85.0\%), reversing the typical
expectation that more examples improve performance.
This suggests that Llama~3.1~8B has already internalised strong
step-by-step reasoning during pre-training, and the 8-shot
demonstrations introduce noise rather than useful signal for this
sample.

ToT (depth=6) matches Few-Shot CoT at 85.0\% but trails Zero-Shot CoT
by $10.0$\,pp; Section~\ref{sec:results:crossmodel} tests whether these
orderings generalise to a stronger backbone.

\subsection{Depth Ablation for Tree-of-Thoughts}
\label{sec:results:depth}

To isolate the effect of search depth, we ran ToT a second time on
the same 100-problem sample with \texttt{depth=3} (branching factor
unchanged at 2), completing all 100 problems.
This mirrors the depth setting used in the original
ToT paper~\cite{yao2023tot} for short-horizon tasks, allowing us to
contrast a shallow search against the deeper \texttt{depth=6}
configuration on identical problems and the same evaluator prompt.
Table~\ref{tab:depth} reports the result, and
Figure~\ref{fig:depth} visualises the accuracy--cost trade-off.

\begin{table}[H]
\centering
\caption{ToT depth ablation on GSM8K (Llama~3.1~8B, BFS-2,
identical evaluator prompt). Accuracies use the corrected GSM8K labels
(see Section~\ref{sec:setup}).}
\label{tab:depth}
\begin{tabular}{@{}lcccc@{}}
\toprule
\textbf{Configuration} & \textbf{Accuracy} & \textbf{N} &
\makecell{\textbf{API calls}\\\textbf{per problem}} &
\makecell{\textbf{Wall-clock}\\\textbf{per problem}} \\
\midrule
ToT BFS-2, depth=3 & 81.0\% & 100 & $\sim$21 & $\sim$3 min \\
ToT BFS-2, depth=6 & 85.0\% & 100 & $\sim$45 & $\sim$6 min \\
\midrule
$\Delta$ (depth=6 $-$ depth=3) & \textbf{+4.0\,pp} & -- &
\textbf{+114\%} & \textbf{+100\%} \\
\bottomrule
\end{tabular}
\end{table}

\begin{figure}[H]
\centering
\includegraphics[width=0.80\textwidth]{depth_ablation}
\caption{ToT depth ablation on GSM8K with Llama~3.1~8B.
Doubling the maximum search depth yields a $+4.0$~pp absolute accuracy
gain while more than doubling the per-problem API call count.}
\label{fig:depth}
\end{figure}

The $+4.0$\,pp gain is small relative to the $2.1\times$ cost increase.
Unlike GPT-4's 100B+ evaluator~\cite{yao2023tot}, our 8B
\textit{sure/maybe/impossible} scorer is weakly discriminative---most
partial thoughts receive \textit{sure} regardless of correctness,
collapsing BFS toward a near-linear chain.
GSM8K problems (2--8 steps) also fit comfortably within depth=3,
so additional depth helps only a small tail of multi-step problems.
The bottleneck at the 8B scale is therefore evaluator discrimination
rather than search depth.

\subsection{Self-Consistency: Vote Agreement Analysis}
\label{sec:results:sc_agreement}

Table~\ref{tab:sc_agreement} breaks down the vote distribution across
the 100 Self-Consistency problems.

\begin{table}[H]
\centering
\caption{Vote agreement distribution for Self-Consistency (5 paths, N=100).}
\label{tab:sc_agreement}
\begin{tabular}{@{}lcc@{}}
\toprule
\textbf{Agreement Level} & \textbf{Count} & \textbf{\%} \\
\midrule
5/5 unanimous         & 65 & 65.0\% \\
4/5 strong majority   & 18 & 18.0\% \\
3/5 bare majority     & 12 & 12.0\% \\
2/5 plurality (no majority) &  3 & 3.0\% \\
1/5 (all distinct)    &  2 & 2.0\% \\
\midrule
Average majority margin & \multicolumn{2}{c}{88.2\%} \\
\bottomrule
\end{tabular}
\end{table}

In 65\% of problems all five paths agreed unanimously; a further 30\%
reached a 4/5 or 3/5 majority, so 95\% of problems admit a clean
Self-Consistency decision.
The 88.2\% average majority margin confirms that 5 paths are sufficient
for reliable consensus on this benchmark.

\subsection{Self-Consistency: Pass@k Analysis}
\label{sec:results:passk}

Table~\ref{tab:passk} reports $\text{Pass@}k$~\cite{chen2021codex}---the
fraction of problems where at least one of the first $k$ paths is correct.

\begin{table}[H]
\centering
\caption{Self-Consistency Pass@$k$ on GSM8K (N=100, gold-corrected).}
\label{tab:passk}
\begin{tabular}{@{}lc@{}}
\toprule
\textbf{Metric} & \textbf{Accuracy} \\
\midrule
Pass@1 (average over 5 samples)   & 86.4\% \\
Pass@3                            & 97.0\% \\
Pass@5 (oracle)                   & \textbf{99.0\%} \\
\midrule
Majority vote @5 (reported)       & 92.0\% \\
Verification gap (oracle $-$ vote) & \textbf{+7.0\,pp} \\
\bottomrule
\end{tabular}
\end{table}

Pass@1 (86.4\%) closely tracks greedy Few-Shot CoT (85.0\%), confirming
that a single stochastic sample adds no consistent gain over the
deterministic baseline.
Pass@5 reaches 99.0\%, exposing a $+7.0$\,pp \textit{verification gap}:
in 7 of 100 problems the correct answer is among the five sampled paths
but loses the majority vote, indicating that a smarter aggregator---not
more samples---is the main opportunity for further gains.

\subsection{Accuracy Stratified by Problem Complexity}
\label{sec:results:stratified}

We adopt the Cobbe et al.~\cite{cobbe2021gsm8k} difficulty split:
problems are bucketed by gold-solution operation count into
\textit{easy} ($<3$), \textit{medium} ($=3$), and \textit{hard}
($>3$) groups (30/33/37 in our sample).
Table~\ref{tab:strat} reports per-bucket accuracy.

\begin{table}[H]
\centering
\caption{Accuracy stratified by problem complexity
(Cobbe et al.~\cite{cobbe2021gsm8k} split, gold-corrected).
N=100, buckets 30/33/37.}
\label{tab:strat}
\begin{tabular}{@{}lccccc@{}}
\toprule
\textbf{Strategy} &
\textbf{Easy ($<$3 ops)} &
\textbf{Medium ($=$3 ops)} &
\textbf{Hard ($>$3 ops)} &
\textbf{Overall} \\
\midrule
Zero-Shot CoT             & \textbf{96.7\%} & \textbf{93.9\%} & \textbf{94.6\%} & \textbf{95.0\%} \\
Few-Shot CoT (8-shot)     & 90.0\% & 87.9\% & 78.4\% & 85.0\% \\
ToT (BFS-2, depth=3)      & 86.7\% & 75.8\% & 81.1\% & 81.0\% \\
ToT (BFS-2, depth=6)      & 90.0\% & 84.8\% & 81.1\% & 85.0\% \\
Self-Consistency (5 paths)& 96.7\% & 93.9\% & 86.5\% & 92.0\% \\
\bottomrule
\end{tabular}
\end{table}

Three results stand out.
First, every strategy shows the expected easy-to-hard accuracy decay.
Second, Few-Shot CoT loses the most on hard problems ($-16.2$\,pp vs.\
Zero-Shot CoT), since its short 2--3-step demonstrations do not
transfer well to longer reasoning chains.
Third, ToT (depth=3) collapses on medium problems (75.8\%) because
3-operation problems exactly exhaust its depth budget, leaving no slack
for BFS recovery after an initial error.
After the format fix, Zero-Shot CoT ties Self-Consistency on easy and
medium buckets and \emph{exceeds} it on hard ($94.6\%$ vs.\ $86.5\%$,
$+8.1$\,pp)---Zero-Shot CoT with disciplined answer extraction is the
strongest strategy at the 8B scale across all difficulty buckets.

\subsection{Solution Length Ratio: Reasoning Verbosity and Overthinking}
\label{sec:results:slr}

Token cost ($\bar{T}$, TPC) does not reveal whether compute is
well-spent.
We quantify reasoning verbosity via the \textbf{Solution Length Ratio}:
\begin{equation}
\text{SLR} = \frac{\text{model reasoning steps}}{\text{gold solution steps}},
\end{equation}
where model steps count \texttt{``= NUMBER''} patterns in the trace
and gold steps are the GSM8K \texttt{<<...>>} annotations.
$\text{SLR}=1$ is canonical length; $\text{SLR}>1$ indicates
overthinking~\cite{wu2025whenmoreisless}.
Table~\ref{tab:slr} reports per-strategy averages.

\begin{table}[H]
\centering
\caption{Solution Length Ratio (SLR): average model reasoning steps
divided by gold solution steps on GSM8K. Higher SLR values indicate
greater overthinking relative to the canonical solution. Average gold
length is 3.3 steps across all strategies (same problem set).}
\label{tab:slr}
\begin{tabular}{@{}lccc@{}}
\toprule
\textbf{Strategy} & \textbf{Avg.\ model steps} &
\textbf{Avg.\ gold steps} & \textbf{SLR} \\
\midrule
Self-Consistency (5 paths)    & 4.2  & 3.3 & \textbf{1.27$\times$} \\
Few-Shot CoT (8-shot)         & 4.1  & 3.3 & 1.44$\times$ \\
Zero-Shot CoT                 & 5.3  & 3.3 & 1.60$\times$ \\
ToT (BFS-2, depth=3)          & 7.9  & 3.3 & 2.72$\times$ \\
ToT (BFS-2, depth=6)          & 9.7  & 3.3 & \textbf{3.30$\times$} \\
\bottomrule
\end{tabular}
\end{table}

Self-Consistency is the most compact strategy (1.27$\times$ SLR), even
below Few-Shot CoT (1.44$\times$), because averaging across five paths
filters out the longest individual chains.
ToT (depth=6), by contrast, inflates solution length to 3.30$\times$
while matching Few-Shot CoT in accuracy---a clean instance of the
\textit{overthinking tax}~\cite{wu2025whenmoreisless}: BFS expansions
add depth without semantic gain because the 8B evaluator cannot prune
unproductive branches.

\subsection{Cross-Model Comparison: Llama~3.1~8B vs.\ Haiku~4.5}
\label{sec:results:crossmodel}

To test whether strategy orderings reflect intrinsic prompting
properties or 8B-scale limitations, we re-ran every strategy on
Claude Haiku~4.5 with full evaluation identity preserved
(Sections~\ref{sec:method:crossmodel}--\ref{sec:setup:crossmodel}).

\paragraph{Headline accuracy.}
Table~\ref{tab:crossmodel} reports overall accuracy on the 100-problem
seed-42 sample under the gold-corrected labels for both models.
\begin{table}[H]
\centering
\caption{Cross-model accuracy (gold-corrected) on the same
100-problem GSM8K sample. ``Cost'' is the realized USD cost for the
full 100-problem run on Haiku~4.5; ZS/FS/SC share a single
700-request batch with the $50\%$ batch discount.}
\label{tab:crossmodel}
\begin{tabular}{@{}lccrc@{}}
\toprule
\textbf{Strategy} & \textbf{Llama~3.1~8B} & \textbf{Haiku~4.5} &
\boldmath$\Delta$ & \textbf{Haiku cost} \\
\midrule
Zero-Shot CoT             & 95.0\% & \textbf{97.0\%} & $+2.0$ &
\multirow{3}{*}{\$0.52 (batch)} \\
Few-Shot CoT (8-shot)     & 85.0\% & \textbf{98.0\%} & $+13.0$ & \\
Self-Consistency (5 paths)& 92.0\% & \textbf{99.0\%} & $+7.0$ & \\
\midrule
ToT (BFS-2, depth=3)      & 81.0\% & \textbf{98.0\%} & $+17.0$ & \$0.72 (sync) \\
ToT (BFS-2, depth=6)      & 85.0\% & \textbf{98.0\%} & $+13.0$ & \$2.02 (sync) \\
\midrule
\multicolumn{4}{r}{\textbf{Total realized cost on Haiku~4.5:}} &
\textbf{\$3.26} \\
\bottomrule
\end{tabular}
\end{table}

The first observation is that every strategy improves on the stronger
backbone, as expected.
The more interesting observation is that the \emph{gaps between
strategies collapse}: on Llama the five strategies span $14$~pp
(81.0\% to 95.0\%), whereas on Haiku they span only $2$~pp
(97.0\% to 99.0\%).
The simplest strategy on Haiku (Zero-Shot CoT, 97.0\%) already
outperforms the most elaborate strategy on Llama
(Self-Consistency, 92.0\%) by $5$~pp---suggesting that at the
benchmark difficulty of GSM8K, raw backbone capability dominates
prompting-strategy choice once the backbone is strong enough.

\paragraph{Stratified view.}
Table~\ref{tab:crossmodel:strat} breaks the Haiku results into the same
Cobbe difficulty buckets (Llama numbers in Table~\ref{tab:strat}).

\begin{table}[H]
\centering
\caption{Cobbe-stratified accuracy on Haiku~4.5
(gold-corrected, same 100-problem sample, buckets 30/33/37).}
\label{tab:crossmodel:strat}
\begin{tabular}{@{}lccccc@{}}
\toprule
\textbf{Strategy} &
\textbf{Easy ($<$3 ops)} &
\textbf{Medium ($=$3 ops)} &
\textbf{Hard ($>$3 ops)} &
\textbf{Overall} \\
\midrule
Zero-Shot CoT             & \textbf{100.0\%} & 97.0\%  & 94.6\%          & 97.0\% \\
Few-Shot CoT (8-shot)     & \textbf{100.0\%} & \textbf{100.0\%} & 94.6\%          & 98.0\% \\
ToT (BFS-2, depth=3)      & \textbf{100.0\%} & \textbf{100.0\%} & 94.6\%          & 98.0\% \\
ToT (BFS-2, depth=6)      & \textbf{100.0\%} & \textbf{100.0\%} & 94.6\%          & 98.0\% \\
Self-Consistency (5 paths)& \textbf{100.0\%} & \textbf{100.0\%} & \textbf{97.3\%} & \textbf{99.0\%} \\
\bottomrule
\end{tabular}
\end{table}

Four of five strategies achieve a perfect 63/63 on Easy and Medium;
Haiku's entire residual error budget lives in the Hard bucket.
Only Self-Consistency recovers one additional Hard answer (36/37 vs.\
35/37), mirroring its 8B advantage---Hard accuracy is the sole
discriminating metric once the backbone nears ceiling.

\paragraph{Cross-model depth ablation: ceiling effect.}
Table~\ref{tab:crossmodel:depth} shows the sharpest single contrast:
depth=3→6 gains $+4.0$~pp on Llama but $0.0$~pp on Haiku at
$2.8\times$ the token cost.
\begin{table}[H]
\centering
\caption{Cross-model ToT depth ablation ($\Delta$ = gold-corrected accuracy).}
\label{tab:crossmodel:depth}
\begin{tabular}{@{}lccc@{}}
\toprule
\textbf{Backbone} &
\textbf{depth=3} & \textbf{depth=6} & \boldmath$\Delta$ \\
\midrule
Llama~3.1~8B   & 81.0\% & 85.0\% & $+4.0$\,pp \\
Haiku~4.5      & 98.0\% & 98.0\% & $\phantom{+}0.0$\,pp \\
\midrule
Haiku token cost &
\multicolumn{3}{r}{depth=6 is $3.1\times$ input and $2.4\times$ output of depth=3} \\
\bottomrule
\end{tabular}
\end{table}
On Haiku neither depth beats Zero-Shot CoT on Hard (both 94.6\%);
ToT's depth parameter buys accuracy only when the base model has
headroom to lose and the evaluator has branches to discriminate.

\paragraph{Self-Consistency verifier saturation.}
On Llama the verifier gap is $+7.0$\,pp (Pass@5 = 99\%,
majority-vote = 92\%): the correct answer is present in the five
paths but discarded by majority voting.
On Haiku the gap is $0.0$\,pp (Pass@5 = majority-vote = 99\%): on
the one remaining error, \emph{none} of the five paths produces the
correct answer.
Path diversity has collapsed (98/100 problems show unanimous 5/5
voting), confirming that Self-Consistency's headroom is a model-scale
artefact, not an intrinsic strategy benefit.

\subsection{Strategy Agreement: Pairwise Cohen's $\kappa$ Across Backbones}
\label{sec:results:kappa}

Accuracy alone does not reveal whether two strategies are correct on
the \emph{same} problems.
We therefore compute all pairwise \textbf{Cohen's $\kappa$}~\cite{cohen1960kappa}
over the same 100-problem gold-corrected sample, treating each strategy
as a binary classifier; $\kappa=0$ is chance, $\kappa=1$ is perfect
agreement (Landis--Koch ranges~\cite{landis1977kappa}).

\begin{table}[H]
\centering
\caption{Pairwise Cohen's $\kappa$ between the five strategies on
Llama~3.1~8B (100-problem seed-42 sample, gold-corrected).}
\label{tab:kappa:llama}
\begin{tabular}{@{}lccccc@{}}
\toprule
& \textbf{ZS-CoT} & \textbf{FS-CoT} & \textbf{SC (5)} &
\textbf{ToT d=3} & \textbf{ToT d=6} \\
\midrule
ZS-CoT v2  & --     & 0.243 & 0.262 & 0.095 & 0.027 \\
FS-CoT     & 0.243  & --    & \textbf{0.563} & 0.435 & 0.216 \\
SC (5)     & 0.262  & \textbf{0.563} & --    & 0.207 & 0.078 \\
ToT d=3    & 0.095  & 0.435 & 0.207 & --    & \textbf{0.011} \\
ToT d=6    & 0.027  & 0.216 & 0.078 & \textbf{0.011} & --    \\
\bottomrule
\end{tabular}
\end{table}

On Llama, ToT depth=3 and depth=6 reach near-zero pairwise agreement
($\kappa=0.011$) despite sharing the same BFS-2 algorithm and evaluator
prompt, confirming that BFS exploration at the 8B scale is stochastic
rather than deliberate.
Zero-Shot CoT shows only fair-to-slight agreement with every other
strategy ($\kappa \in [0.027, 0.262]$), reflecting the genuine
cross-strategy complementarity that drives the $+7.0$\,pp Pass@5
headroom (Section~\ref{sec:results:passk}).

\begin{table}[H]
\centering
\caption{Pairwise Cohen's $\kappa$ between the five strategies on
Claude Haiku~4.5 (same 100-problem sample, gold-corrected).}
\label{tab:kappa:haiku}
\begin{tabular}{@{}lccccc@{}}
\toprule
& \textbf{ZS-CoT} & \textbf{FS-CoT} & \textbf{SC (5)} &
\textbf{ToT d=3} & \textbf{ToT d=6} \\
\midrule
ZS-CoT v2  & --     & 0.385 & 0.492 & \textbf{0.795} & \textbf{0.795} \\
FS-CoT     & 0.385  & --    & 0.662 & 0.490 & 0.490 \\
SC (5)     & 0.492  & 0.662 & --    & 0.662 & 0.662 \\
ToT d=3    & \textbf{0.795} & 0.490 & 0.662 & --    & \textbf{1.000} \\
ToT d=6    & \textbf{0.795} & 0.490 & 0.662 & \textbf{1.000} & --    \\
\bottomrule
\end{tabular}
\end{table}

On Haiku the picture inverts: ToT depth=3 and depth=6 reach
$\kappa=1.000$---perfect problem-by-problem concordance---placing the
ceiling effect on a per-problem rather than only aggregate basis.
The overall $\kappa$ range collapses from $[0.011,\,0.563]$ on Llama
to $[0.385,\,1.000]$ on Haiku: \emph{strategy independence itself
decays with model scale}.
ToT (either depth) agrees with Zero-Shot CoT at $\kappa=0.795$,
indicating that deliberate-search machinery collapses into single-chain
behaviour once the backbone is strong enough.

\subsection{Ablation Study}

\paragraph{Temperature sensitivity.}
Greedy decoding (88.0\%) outperforms stochastic sampling at $T{=}0.7$
(86.0\%) by 2\,pp on 50 samples, confirming a marginal but consistent
advantage for deterministic decoding on arithmetic tasks
(Table~\ref{tab:temp}).

\begin{table}[H]
\centering
\caption{Temperature ablation, Few-Shot CoT (N=50, seed=42).}
\label{tab:temp}
\begin{tabular}{@{}lcc@{}}
\toprule
\textbf{Temperature} & \textbf{Accuracy} & \textbf{N} \\
\midrule
0.0 (greedy)     & \textbf{88.0\%} & 50 \\
0.7 (stochastic) & 86.0\%          & 50 \\
\bottomrule
\end{tabular}
\end{table}

\paragraph{Prompt sensitivity.}
Accuracy across three 8-shot example sets spans 6\,pp (82.0\%--88.0\%),
indicating moderate sensitivity to example choice (Table~\ref{tab:prompt}).
Set-1 underperforms (82.0\%), while the canonical Wei et al.\ examples
and Set-2 both reach 88.0\%.

\begin{table}[H]
\centering
\caption{Prompt sensitivity, Few-Shot CoT (temperature=0, N=50, seed=42).}
\label{tab:prompt}
\begin{tabular}{@{}lcc@{}}
\toprule
\textbf{Example Set} & \textbf{Accuracy} & \textbf{N} \\
\midrule
Set-0 (canonical, Wei et al.) & \textbf{88.0\%} & 50 \\
Set-1 (random, seed=1)        & 82.0\%          & 50 \\
Set-2 (random, seed=2)        & \textbf{88.0\%} & 50 \\
\bottomrule
\end{tabular}
\end{table}

%------------------------------------------------------------------------
\section{Discussion}
\label{sec:discussion}

\medskip\noindent\textbf{ToT matches but does not beat CoT on 8B models.}

ToT (depth=6) achieves 85.0\% on 100 samples, closely matching the CoT
baselines but not outperforming them as Yao et
al.~\cite{yao2023tot} report on GPT-4.
The depth ablation in Section~\ref{sec:results:depth} reinforces
this picture: increasing depth from 3 to 6 yields only $+4.0$\,pp at
$2.1\times$ the cost.
Two factors plausibly account for the absence of a gain.
First, GSM8K problems are typically solvable in 3--5 chained steps,
which CoT already handles well---the multi-branch exploration ToT
offers carries diminishing returns on this benchmark.
Second, the evaluator (\textit{sure/maybe/impossible}) is implemented
with the same 8B model: at this scale, the evaluator's discriminative
power is weak, which collapses BFS toward a near-linear search rather
than a true tree exploration.

\medskip\noindent\textbf{Zero-shot CoT substantially outperforms Few-shot CoT.}

The 10-point accuracy gap between Zero-Shot CoT (95.0\%) and
Few-Shot CoT (85.0\%) suggests that Llama~3.1~8B's pre-training has
already internalised the step-by-step reasoning style, and the
8-shot demonstrations from Wei et al.\ introduce systematic noise
rather than useful signal.
This goes further than the trend observed in recent
instruction-tuned models, where zero-shot and few-shot performance
often only \emph{converge}: here Zero-Shot is a strict $+10$\,pp
better, indicating that the prompt-format mismatch (the demos
encourage ``\textit{The answer is X.}'' while the model's free-form
chains carry richer reasoning) is actively harmful.

\medskip\noindent\textbf{Self-Consistency no longer dominates after the format fix.}

With the format-aligned Zero-Shot CoT baseline at 95.0\%, Self-Consistency
(92.0\%, at 5$\times$ the per-problem cost) is strictly Pareto-dominated.
The 5-path ensemble still helps on a per-path basis (Pass@1 of $86.4\%$
is below Zero-Shot CoT's 95.0\%; the majority vote partially recovers
to 92.0\%), but the vote aggregation cannot make up the ground gained
by a single well-extracted greedy chain at this model scale.

\medskip\noindent\textbf{Accuracy--cost trade-off.}

Figure~\ref{fig:scatter} visualises the accuracy--cost trade-off across
all five strategy configurations on a log-scale cost axis.
Zero-Shot CoT ($\sim$1 API call, 95.0\%) lies alone on the Pareto
frontier; all other configurations are dominated.
Self-Consistency costs $5\times$ Zero-Shot CoT for $-3.0$\,pp
accuracy; Few-Shot CoT matches the cost at $-10.0$\,pp; ToT (depth=3)
costs $21\times$ for $-14.0$\,pp; and ToT (depth=6) costs $45\times$
for $-10.0$\,pp.
For deployments where accuracy is critical and inference cost is secondary,
Self-Consistency is the optimal choice; for cost-sensitive applications,
Zero-Shot CoT provides the best accuracy-per-call ratio.

\begin{figure}[H]
\centering
\includegraphics[width=0.95\textwidth]{scatter_accuracy_cost}
\caption{Accuracy--cost trade-off across all five prompting
configurations (Llama~3.1~8B, GSM8K).
The horizontal axis is the average number of API calls per problem
(log scale); the dashed line marks the Pareto frontier.
Zero-Shot CoT is Pareto-efficient; Few-Shot CoT, Self-Consistency,
and both Tree-of-Thoughts configurations are dominated.}
\label{fig:scatter}
\end{figure}

\medskip\noindent\textbf{Generalisation across model scales.}

The Haiku~4.5 replication reproduces three qualitative findings.
First, Zero-Shot CoT remains competitive with Few-Shot CoT (97\% vs.\
98\%): demonstrations no longer hurt at the stronger scale, consistent
with the noise argument.
Second, Self-Consistency retains a scale-stable Hard-bucket advantage
(97.3\% on Haiku, 86.5\% on Llama).
Third, ToT (depth=6) matches but does not beat CoT on either
backbone---the absence of a ToT advantage is not an artefact of an
undersized evaluator.
The one reversal: strategy spread collapses from $14$~pp (Llama) to
$2$~pp (Haiku), demonstrating that raw backbone capability dominates
prompting-strategy choice at sufficient scale.

\medskip\noindent\textbf{ToT depth saturates as the base model strengthens.}

Llama gains $+4.0$\,pp going from depth=3 to depth=6; Haiku gains
$0.0$\,pp at $2.8\times$ the token cost.
When the baseline already approaches ceiling, the Easy/Medium tail is
saturated (Table~\ref{tab:crossmodel:strat}) and residual Hard-bucket
errors are arithmetic slips that BFS cannot prune.
Self-Consistency tells the same story from the ensemble side: on
Haiku none of the five paths recovers the problems majority voting
gets wrong, so an oracle verifier cannot help either.
The pairwise $\kappa$ analysis (Section~\ref{sec:results:kappa})
independently confirms this: Llama's ToT depths reach $\kappa=0.011$
(stochastic exploration) while Haiku's reach $\kappa=1.000$
(deterministic ceiling).

\medskip\noindent\textbf{Error analysis and limitations.}

Manual inspection of incorrect predictions reveals that arithmetic
errors account for the vast majority of failures across all strategies:
models set up the correct reasoning structure but miscompute at one
step, with the error propagating to the final answer.
On Haiku~4.5 the same pattern holds at a much smaller budget---all
residual errors are arithmetic rather than misinterpretation or step
omission, confirming that arithmetic brittleness persists even at
frontier scale, only attenuated.
Beyond this, the main limitations are: (1)~exact-match accuracy ignores
reasoning quality; (2)~the cross-model comparison covers only two
families---adding Mistral, Qwen, or Gemma at the 7--8B scale would
separate pre-training data quality from raw scale; (3)~GSM8K's short
arithmetic horizon means our findings should not be extrapolated to
long-horizon planning benchmarks where ToT reported larger
gains~\cite{yao2023tot}.

%------------------------------------------------------------------------
\section{Conclusion}
\label{sec:conclusion}

We evaluated Zero-Shot CoT, Few-Shot CoT, ToT (at two search depths),
and Self-Consistency on GSM8K using Llama~3.1~8B.
Zero-Shot CoT (95.0\%, after the answer-extraction format fix) is both
the most accurate and the most cost-efficient strategy---a strict
Pareto dominator over all other configurations.
ToT (depth=6) matches but does not beat Few-Shot CoT at 85.0\%;
a depth ablation shows that halving the search depth to 3 costs only
$4.0$\,pp while reducing inference cost by more than half, indicating
diminishing returns at this scale due to weak evaluator discrimination.

A cross-model replication on Claude Haiku~4.5 with exact sample,
prompt, parser, and label-correction identity preserved collapses the
strategy gap from $14$~pp to $2$~pp.
The depth ablation contrast is sharper still: on the stronger
backbone, doubling ToT's search depth yields exactly zero additional
accuracy at $2.8\times$ the token cost, and the Self-Consistency
verifier gap shrinks from $+7.0$~pp on Llama to $0.0$~pp on Haiku.
Together, these results point to a \emph{ceiling effect}: deliberation-
and ensemble-style strategies buy real accuracy only when the base
model has headroom to lose and identifiable wrong branches to
prune---conditions that disappear as the backbone strengthens.

\paragraph{Future work.}
Replacing the categorical \textit{sure/maybe/impossible} evaluator
with a learned scorer is the most direct route to better ToT
performance on 8B-class models.
The $7.0$\,pp verification gap between Pass@5 (99.0\%) and majority
vote on Llama further suggests that an external verifier---rather
than larger ensembles---is the main opportunity for gains at fixed
sample budget.
Extending the comparison to additional open-weight families
(e.g.\ Mistral, Qwen, Gemma at the 7--8B scale) would disentangle
pre-training data quality from raw parameter count, and test whether
the ceiling effect generalises across providers.

%------------------------------------------------------------------------
\bibliographystyle{splncs04}
\begin{thebibliography}{99}

\bibitem{brown2020gpt3}
Brown, T., et al.:
Language models are few-shot learners.
In: Advances in Neural Information Processing Systems (NeurIPS), vol. 33,
pp. 1877--1901 (2020)

\bibitem{cobbe2021gsm8k}
Cobbe, K., et al.:
Training verifiers to solve math word problems.
arXiv preprint arXiv:2110.14168 (2021)

\bibitem{kojima2022zeroshot}
Kojima, T., Gu, S.S., Reid, M., Matsuo, Y., Iwasawa, Y.:
Large language models are zero-shot reasoners.
In: Advances in Neural Information Processing Systems (NeurIPS), vol. 35,
pp. 22199--22213 (2022)

\bibitem{wang2022selfconsistency}
Wang, X., et al.:
Self-consistency improves chain of thought reasoning in language models.
In: International Conference on Learning Representations (ICLR) (2023)

\bibitem{wei2022cot}
Wei, J., et al.:
Chain-of-thought prompting elicits reasoning in large language models.
In: Advances in Neural Information Processing Systems (NeurIPS), vol. 35,
pp. 24824--24837 (2022)

\bibitem{yao2023tot}
Yao, S., et al.:
Tree of thoughts: Deliberate problem solving with large language models.
In: Advances in Neural Information Processing Systems (NeurIPS), vol. 36
(2023)

\bibitem{chen2021codex}
Chen, M., et al.:
Evaluating large language models trained on code.
arXiv preprint arXiv:2107.03374 (2021)

\bibitem{ockbench2025}
OckBench Authors:
OckBench: Measuring the efficiency of LLM reasoning.
arXiv preprint arXiv:2511.05722 (2025)

\bibitem{tale2024}
Han, T., et al.:
Token-budget-aware LLM reasoning.
arXiv preprint arXiv:2412.18547 (2024)

\bibitem{zhang2022autocot}
Zhang, Z., Zhang, A., Li, M., Smola, A.:
Automatic chain of thought prompting in large language models.
In: International Conference on Learning Representations (ICLR) (2023)

\bibitem{zhou2022ltm}
Zhou, D., et al.:
Least-to-most prompting enables complex reasoning in large language models.
In: International Conference on Learning Representations (ICLR) (2023)

\bibitem{wu2025whenmoreisless}
Wu, Y., et al.:
When more is less: Understanding chain-of-thought length in LLMs.
arXiv preprint arXiv:2502.07266 (2025)

\bibitem{cohen1960kappa}
Cohen, J.:
A coefficient of agreement for nominal scales.
Educational and Psychological Measurement \textbf{20}(1), 37--46 (1960)

\bibitem{landis1977kappa}
Landis, J.R., Koch, G.G.:
The measurement of observer agreement for categorical data.
Biometrics \textbf{33}(1), 159--174 (1977)

\bibitem{artstein2008agreement}
Artstein, R., Poesio, M.:
Inter-coder agreement for computational linguistics.
Computational Linguistics \textbf{34}(4), 555--596 (2008)

\end{thebibliography}

\end{document}

